\subsection{Nagata's theorem}

THEOREM 1 (Tate). --- Let A be a Noetherian integrally closed ring, and a an element of A. Suppose that the ideal aA is prime, that the ring A/aA is Japanese, and that A is complete for the aA-adic topology. Then the ring A is Japanese.

\begin{enumerate}
\def\labelenumi{\alph{enumi})}
\tightlist
\item
  Let K be the field of fractions of A. The assertion being trivial when K has characteristic 0 (no. 1, corollary of prop. 1), one may suppose K has characteristic \(p > 0\) . One may also suppose \(a \neq 0\) .
\end{enumerate}

Let L be a finite radical extension of K and q a power of p such that \(L \subset K^{1/q}\) . Put \(x = a^{1/q}\) and \(M = L(x)\) . By prop. 1 of no. 1, it suffices to prove that the integral closure B of A in M is a finite A-algebra.

\begin{enumerate}
\def\labelenumi{\alph{enumi})}
\setcounter{enumi}{1}
\item
  Let us first prove that the ideal \(xB\) is the unique prime ideal of B lying over \(aA\) . Indeed, there exists at least one prime ideal of B lying over \(aA\) (V, § 2, no. 1, Th. 1). Let \(\mathfrak{q}\) be one of these ideals. We have \(x^{q} = a \in \mathfrak{q}\) , whence \(xB \subset \mathfrak{q}\) since \(\mathfrak{q}\) is prime. Conversely, let \(y\) be an element of \(\mathfrak{q}\); the element \(y^{q}\) of K is integral over A, hence belongs to A since A is integrally closed. Since \(\mathfrak{q} \cap A = aA\) , there exists an element \(\alpha\) of A such that \(y^{q} = a\alpha = x^{q}\alpha\) . Consequently, the element \(y/x\) of M is integral over A, hence belongs to B; thus we have \(y \in xB\) , whence \(\mathfrak{q} = xB\) , which proves our assertion.
\item
  It follows that the ring \(B_{xB}\) is the integral closure in M of the ring \(A_{aA}\) (V, § 1, no. 5, prop. 16 and § 2, no. 1, prop. 2). By VI, § 3, no. 6, prop. 9, \(A_{aA}\) is a discrete valuation ring; one then deduces from the Krull-Akizuki theorem (VII, § 2, no. 5, prop. 5) that the field \(\kappa(xB)\) is a finite extension of \(\kappa(aA)\) and that \(B_{xB}\) is Noetherian.
\item
  The ring B/xB is integral over the Japanese ring A/aA and its field of fractions is a finite extension of the field of fractions of the latter. Consequently, B/xB is a finitely generated (A/aA)-module. For every integer \(i \geqslant 0\) , the same is true of the module \(x^{i}B/x^{i+1}B\) . Consequently, the (A/aA)-module B/aB has a composition series of length q whose quotients are finitely generated (A/aA)-modules, and hence B/aB is itself a finitely generated (A/aA)-module.
\item
  Equip the ring A with the (aA)-adic filtration and the ring B with the (aB)-adic filtration. Then A is complete by hypothesis; since \(B_{xB}\) is integral and Noetherian, the \(aB_{xB}\)-adic filtration of \(B_{xB}\) is separated (III, § 3, no. 2, corollary to prop. 5); consequently we have \(\bigcap a^{n}B \subset \bigcap a^{n}B_{xB} = \{0\}\) and the \(aB\)-adic filtration of B is separated; the gr(A)-module gr(B) is generated by gr₀(B), hence is finitely generated by d). It then follows from III, § 2, no. 9, cor. 1 to prop. 12, that B is a finitely generated A-module, which completes the proof.
\end{enumerate}

COROLLARY. --- Let R be a Noetherian integral domain and n an integer. If R is Japanese, the ring \(R[[T_{1}, \ldots, T_{n}]]\) is Japanese.

Reasoning by induction, we may assume \(n = 1\) . Let S denote the integral closure of R; if R is Japanese, S is a finite algebra over R, hence a Japanese ring (no. 1, prop. 2). The ring \(S[[T]]\) is Noetherian and integrally closed (V, § 1, no. 4, prop. 14); applying Thm. 1 to A = \(S[[T]]\) and \(a = T\) , we deduce that \(S[[T]]\) is Japanese. Consequently \(R[[T]]\) is Japanese (no. 1, prop. 2).

THEOREM 2 (Nagata). --- Every complete local Noetherian integral domain A is Japanese.

By Th. 3 of § 2, no. 5 and Th. 2 of § 3, no. 3, there exist an integer \(n \geqslant 0\) , a ring R that is either a field or a discrete valuation ring whose field of fractions has characteristic 0, and a subring B of A, isomorphic to \(R[[T_{1}, \ldots, T_{n}]]\), such that A is a finite B-algebra. Then R is Japanese (no. 1, example and corollary of prop. 1), hence B is Japanese (corollary to Th. 1), and A is Japanese (no. 1, prop. 2).

COROLLARY. --- Let A be a Noetherian semi-local ring whose completion is reduced. Then the integral closure A' of A in its total ring of fractions R is a finite A-algebra.

Suppose first that A is local and complete, and let \(p_{1}, \ldots, p_{n}\) be the distinct minimal prime ideals of A; for \(i = 1, \ldots, n\) , denote by \(K_{i}\) the field of fractions of \(A/p_{i}\) and by \(A_{i}'\) the integral closure of \(A/p_{i}\) . Since A is reduced, R is the product of the rings \(K_{i}\) and \(A'\) is the product of the rings \(A_{i}'\) (V, § 1, no. 2, cor. 1 to prop. 9). Since the local rings \(A/p_{i}\) are integral and complete, they are Japanese (Th. 2), so that each \(A_{i}'\) is a finite A-algebra, and \(A'\) is a finite A-algebra.

If A is semi-local and complete, it is isomorphic to a finite product of complete local rings (III, § 2, no. 13, corollary to prop. 19), and one concludes immediately from what precedes.

Let us turn to the general case and note that the completion \(\hat{\mathbf{A}}\) of A is a semi-local, complete, Noetherian ring, faithfully flat over A (III, loc. cit., § 3, no. 4, corollary of prop. 8 and § 3, no. 5, prop. 9). Let S be the set of non-zero-divisors of A; we have R = S \(^{-1}\) A. Since \(\hat{\mathbf{A}}\) is flat over A, the elements of S are non-zero-divisors in \(\hat{\mathbf{A}}\) and S \(^{-1}\) \(\hat{\mathbf{A}}\) is identified with a subring of the total ring of fractions T of \(\hat{\mathbf{A}}\) . Still since \(\hat{\mathbf{A}}\) is flat over A, the ring A′ ⊗ \(_{\mathbf{A}}\) \(\hat{\mathbf{A}}\) is identified with a subring of R ⊗ \(_{\mathbf{A}}\) \(\hat{\mathbf{A}}\) = S \(^{-1}\) \(\hat{\mathbf{A}}\) , hence also with a subring of T integral over \(\hat{\mathbf{A}}\) . By the first part of the proof, A′ ⊗ \(_{\mathbf{A}}\) \(\hat{\mathbf{A}}\) is therefore a \(\hat{\mathbf{A}}\) -module of finite type; consequently, A′ is a finitely generated A-module (I, § 3, no. 6, prop. 11).

Recall (A, V, p. 114, def. 1) that an algebra E over a field K is said to be separable if the ring L \(\otimes_{\mathbf{K}}\) E is reduced for every extension L of K; it suffices that this be the case for every finite extension of K. The following proposition generalizes Th. 2:

PROPOSITION 3. --- Let A be a Noetherian semi-local integral domain, and K its field of fractions. If the K-algebra \(K \otimes_{A} \hat{\mathbf{A}}\) is separable, the ring A is Japanese.

Let L be a finite extension of K and let B be the integral closure of A in L. Let F be a finite subset of B such that \(L = K[F]\) (V, § 1, no. 5, cor. 2 to prop. 16); let C denote the (finite) A-algebra generated by F. Since L is the field of fractions of C, the ring B is the integral closure of C (V, § 1, no. 1, prop. 6), and it suffices to prove that B is a finite C-algebra. Now, C is a semi-local Noetherian ring (IV, § 2, no. 5, cor. 3 to prop. 9); its completion is identified with \(C \otimes_{A} \hat{A}\) (III, § 3, no. 4, th. 3 (ii)), hence also to a subring of the reduced ring \(L \otimes_{A} \hat{A} = L \otimes_{K} (K \otimes_{A} \hat{A})\) and consequently is reduced. Prop. 3 therefore follows from the corollary to Th. 2.

